% ====================================================================
% FILE: content/11_operators_full.tex
% Полное семейство операторов F, G, H, Q, R, S, U, Psi, E_int
% Core 1.3.3 — canonical version
% ====================================================================

\section{Операторы эволюции и взаимодействия}
\label{sec:operators-full}

Эта глава разворачивает абстрактный оператор эволюции, введённый в
разделе~the formal model chapter, в полное семейство операторов, используемое в
онтологии континуумов (OC). Она объединяет определения, разработанные в
archival source corpus, и представляет их в компактной, вертикально согласованной форме,
пригодной для Core~1.3.3.

На протяжении этой секции континуум записывается как
\[
  K(t) =
  \big(
    \Omega(t), A(t), P(t), J(t),
    \Theta(t), \partial\Omega(t), C(t), k(t)
  \big),
\]
с компонентами, определёнными в разделе~the formal model chapter. Пространство
внедрения обозначается через \(M\), с осями \(A(M)\) и ограничениями,
которые задают допустимые состояния \(K\).

\subsection{Декомпозиция оператора эволюции}

Структурная эволюция континуума за бесконечно малый шаг времени
описана оператором
\[
  E : K(t) \longrightarrow K(t+dt).
\]
Вместо указания одного монолитного отображения OC факторизует \(E\) в семейство
операторов, действующих на различные компоненты континуума:
\[
  E = (F, G, H, Q, R, S, U),
\]
где
\begin{itemize}
  \item \(F\) управляет динамикой потоков \(J\);
  \item \(G\) управляет динамикой потенциалов \(P\);
  \item \(H\) управляет динамикой порогов \(\Theta\);
  \item \(Q\) управляет динамикой циклов \(C\);
  \item \(R\) управляет динамикой границы \(\partial\Omega\);
  \item \(S\) управляет структурной перенастройкой осей \(A\) и внутренней
        связности;
  \item \(U\) управляет эволюцией континуальности \(k(t)\).
\end{itemize}

На формальном уровне это можно записать как
\begin{align*}
  J(t+dt)              &= F\big(J(t), P(t), A(t), \Theta(t), M\big), \\
  P(t+dt)              &= G\big(P(t), A(t), J(t), \Theta(t), M\big), \\
  \Theta(t+dt)         &= H\big(\Theta(t), P(t), J(t), M\big), \\
  C(t+dt)              &= Q\big(C(t), P(t), J(t), \Theta(t)\big), \\
  \partial\Omega(t+dt) &= R\big(\partial\Omega(t), P(t), J(t), \Theta(t), M\big), \\
  A(t+dt)              &= S\big(A(t), P(t), J(t), \Theta(t), M\big), \\
  k(t+dt)              &= U\big(k(t), \Omega(t), A(t), P(t), J(t),
                               \Theta(t), C(t), \partial\Omega(t), M\big).
\end{align*}

Операторы ограничиваются аксиомами OC, в частности монотонностью размерности,
невозможностью эволюции вне пространства внедрения и определением рождения и
смерти.

\subsection{Оператор \texorpdfstring{\texorpdfstring{$F$}{F}}{F}: динамика потоков}

Оператор потоков \(F\) определяет временное изменение потоков:
\[
  J(t+dt) = F\big(J(t), P(t), A(t), \Theta(t), M\big).
\]

Интуитивно \(F\) кодирует:
\begin{itemize}
  \item законы сохранения или баланса (например, уравнения непрерывности,
        суммы Кирхгофа, сохранение вероятности);
  \item конститутивные соотношения, связывающие потоки с градиентами потенциалов
        (например, законы Фика, Ома, Фурье, кинетику реакций);
  \item структурную классификацию потоков на поддерживающие, критические и
        деструктивные компоненты.
\end{itemize}

Общая декомпозиция имеет вид
\[
  J(t) = J_{\mathrm{support}}(t)
       + J_{\mathrm{critical}}(t)
       + J_{\mathrm{kill}}(t),
\]
а \(F\) сохраняет эту структуру:
\begin{align*}
  J_{\mathrm{support}}(t+dt)
    &= F_{\mathrm{support}}\big(J, P, A, \Theta, M\big), \\
  J_{\mathrm{critical}}(t+dt)
    &= F_{\mathrm{critical}}\big(J, P, A, \Theta, M\big), \\
  J_{\mathrm{kill}}(t+dt)
    &= F_{\mathrm{kill}}\big(J, P, A, \Theta, M\big).
\end{align*}

На разных уровнях \(K_x\) эти классы имеют разные реализации:
\begin{itemize}
  \item в \(K_2\) (физические континуумы) потоки включают диффузионные,
        конвекционные и переносимые полем токи;
  \item в \(K_3\) (химические континуумы) потоки соответствуют скоростям
        реакций и транспорта;
  \item в \(K_4\)–\(K_5\) (протоклеточные и биоэлектрические
        континуумы) потоки включают ионные токи, осмотические потоки и
        метаболические потоки;
  \item в \(K_6\)–\(K_8\) (когнитивные и социальные континуумы) потоки
        становятся потоками информации, внимания, ресурсов и обязательств.
\end{itemize}

\subsection{Оператор \texorpdfstring{\texorpdfstring{$G$}{G}}{G}: динамика потенциалов}

Оператор потенциалов \(G\) определяет, как меняются потенциалы:
\[
  P(t+dt) = G\big(P(t), A(t), J(t), \Theta(t), M\big).
\]

Во многих конкретных моделях \(G\) выражается дифференциальным соотношением
\[
  \frac{dP}{dt} = J,
\]
возможны дополнительные нелинейные слагаемые для диссипации, приведения в
движение и связи. На структурном уровне OC требует лишь, чтобы:
\begin{itemize}
  \item потенциалы реагировали на потоки и, в свою очередь, могли их
        ограничивать;
  \item изменения потенциалов учитывали знак и неравенственную структуру,
        задаваемую порогами \(\Theta\);
  \item при отсутствии потоков \(J = 0\) потенциалы стремились к пороговым
        поверхностям или релаксировали к внутренним предпочтительным
        конфигурациям.
\end{itemize}

Различные классы потенциалов (энергетические, химические, электрокхимические,
информационные, нормативные) являются экземплярами единой структурной роли:
закодирования ограничений и движущих сил на \(\Omega(K)\).

\subsection{Оператор \texorpdfstring{\texorpdfstring{$H$}{H}}{H}: динамика порогов}

Пороги не статичны; они могут адаптироваться, дрейфовать или
реконфигурироваться под влиянием потенциалов, потоков и ограничений
внедрения. Оператор порогов \(H\) определён следующим образом:
\[
  \Theta(t+dt) = H\big(\Theta(t), P(t), J(t), M\big).
\]

К таким примерам относятся:
\begin{itemize}
  \item адаптация кривых активации ионных каналов в нервной ткани;
  \item пластичность диапазонов жизнеспособности для протоклеток при изменяющейся
        среде;
  \item сдвиг норм и институциональных порогов в социальных системах;
  \item изменения порогов когерентности или непротиворечивости в
        теоретических системах.
\end{itemize}

Формально \(H\) должен удовлетворять:
\begin{itemize}
  \item совместимости с ограничениями внедрения: изменения порогов не могут
        требовать конфигураций, запрещённых \(M\);
  \item сохранению таксономии порогов
        (\(\Theta_{\mathrm{exist}}, \Theta_{\mathrm{stab}},
        \Theta_{\mathrm{crit}}, \Theta_{\mathrm{dim}},
        \Theta_{\mathrm{death}}, \Theta_{\mathrm{expr}}, \Theta_{\mathrm{embed}}\));
  \item локальной непрерывности: малые изменения потенциалов или потоков могут
        вызывать малые изменения порогов, за исключением критических точек,
        где допускаются бифуркации.
\end{itemize}

\subsection{Оператор \texorpdfstring{\texorpdfstring{$Q$}{Q}}{Q}: динамика циклов}

Циклы \(C(t)\) представляют собой замкнутые траектории, поддерживающие
организацию континуума. Их эволюция задаётся
\[
  C(t+dt) = Q\big(C(t), P(t), J(t), \Theta(t)\big).
\]

Структурно \(Q\) учитывает:
\begin{itemize}
  \item \emph{рождение} новых циклов, когда потоки замыкаются в пространстве
        состояний вдали от \(\partial\Omega\);
  \item \emph{стабилизацию} или \emph{укрепление} существующих циклов, когда
        доминируют поддерживающие потоки;
  \item \emph{ослабление} и \emph{разрушение} циклов под действием
        деструктивных потоков или нарушения порогов.
\end{itemize}

Особое значение имеют комплексы циклов. Пусть \(C_{\max}(t)\) обозначает
максимальное множество взаимно совместимых циклов, поддерживающих континуум.
Тогда в теореме~9 раздела~the section on resultsутверждается, что смерть
совпадает с условием \(C_{\max}(t) = \emptyset\). Поэтому оператор \(Q\)
непосредственно участвует в переходах жизнь–смерть.

\subsection{Оператор \texorpdfstring{\texorpdfstring{$R$}{R}}{R}: эволюция границы}

Оператор границы \(R\) управляет геометрией и положением границы
\(\partial\Omega\):
\[
  \partial\Omega(t+dt) =
    R\big(\partial\Omega(t), P(t), J(t), \Theta(t), M\big).
\]

На структурном уровне \(R\) должен:
\begin{itemize}
  \item быть согласованным с определениями порогов:
        \(\partial\Omega = \{ s \mid \exists k: f_k(s) = 0 \}\);
  \item допускать расширение, сжатие и бифуркацию \(\Omega\);
  \item кодировать влияние потоков и смещений порогов на доступные регионы.
\end{itemize}

В конкретных реализациях \(R\) может принимать различные формы:
\begin{itemize}
  \item в \(K_2\) он может соответствовать движению фазовых границ,
        перколяционных порогов или поверхностей когерентности;
  \item в \(K_3\)–\(K_4\) он включает динамику мембран, интерфейсов и
        границ компартментов;
  \item в \(K_5\) соответствует порогам возбудимости и рефрактерным
        областям в пространстве проводимости;
  \item в \(K_7\)–\(K_8\) он может представлять институциональные и
        юридические границы.
\end{itemize}

Раздел~the boundary discussionextendedдаёт расширенный разбор геометрии
границы и модели «пластинок».

\subsection{Оператор \texorpdfstring{\texorpdfstring{$S$}{S}}{S}: структурная
перенастройка}

Структурный оператор \(S\) управляет изменениями осей и внутренней
конфигурации:
\[
  A(t+dt) = S\big(A(t), P(t), J(t), \Theta(t), M\big).
\]

Типичные роли \(S\) включают:
\begin{itemize}
  \item активацию или деактивацию осей (например, включение латентных степеней
        свободы);
  \item перестройку паттернов связности внутри континуума;
  \item укрупнение масштаба или уточнение эффективных осей в ходе
        ренормализации или абстракции.
\end{itemize}

Ключевым образом \(S\) подчиняется монотонности размерности:
\[
  \dim A(t+dt) \ge \dim A(t).
\]
Если добавляется новая ось \(A_{\mathrm{new}}\), такая что
\(A_{\mathrm{new}} \notin \mathrm{span}_{\mathrm{str}}(A(t))\), то
происходит переход размерности, и континуум переходит в \(K_{x+1}\), а не
остаётся на том же уровне. Это формализуется через операторы рождения в
разделе~the birth discussionoperatorsниже.

\subsection{Оператор \texorpdfstring{\texorpdfstring{$U$}{U}}{U}: динамика
континуальности}

Оператор \(U\) управляет эволюцией континуальности \(k(t)\):
\[
  k(t+dt) =
  U\big(k(t), \Omega(t), A(t), P(t), J(t),
        \Theta(t), C(t), \partial\Omega(t), M\big).
\]

Используя унифицированное определение \(k(t)\) из раздела~the formal model chapter,
\(U\) оценивает, как изменения пространства состояний, осей, потоков,
порогов, циклов и границы влияют на жизнеспособность континуума. Качественно:
\begin{itemize}
  \item поддерживающие потоки и устойчивые циклы увеличивают или поддерживают
        \(k(t)\);
  \item деструктивные потоки и потеря циклов уменьшают \(k(t)\);
  \item приближение к \(\partial\Omega\) снижает \(k(t)\);
  \item расширение \(\Omega\) и обогащение комплексов циклов повышают
        \(k(t)\).
\end{itemize}

Смерть соответствует \(k(t) \to 0\) вместе с \(\Omega = \emptyset\).
Рождение соответствует появлению нового континуума с \(k(t) > 0\) в
ранее пустой области пространства возможностей.

\subsection{Операторы рождения \texorpdfstring{$\Psi_{x\to x+1}$}{Psi\_x→x+1}
\label{sec:birth-operators}

Размерностные переходы между уровнями опосредуются операторами рождения
\[
  \Psi_{x\to x+1} : (K_x, M_x) \longrightarrow (K_{x+1}, M_{x+1}).
\]

Структурно \(\Psi_{x\to x+1}\) определён всякий раз, когда выполняются
следующие условия:
\begin{enumerate}
  \item \textbf{Новые различия.} \
        Существует класс различий вне текущего структурного
        охвата:
        \[
          D_{\mathrm{new}}\not\subseteq
          \mathrm{span}_{\mathrm{str}}(A(K_x)).
        \]
  \item \textbf{Напряжение выше порога размерности.} \
        Структурное напряжение превышает порог размерности:
        \[
          T(K_x,t) > \Theta_{\mathrm{dim}}(K_x).
        \]
  \item \textbf{Доступная ось в пространстве внедрения.} \
        Пространство внедрения \(M_x\) содержит хотя бы одну ось, подходящую для
        размещения новых различий:
        \[
          A_{\mathrm{new}} \in A(M_x) \setminus A(K_x).
        \]
  \item \textbf{Непустой допустимый регион.} \
        Существует непустое множество состояний \(\Omega(K_{x+1})\), совместимое
        с новой осью и порогами.
\end{enumerate}

Действие \(\Psi_{x\to x+1}\) можно суммарно записать как
\begin{align*}
  A(K_{x+1})      &= A(K_x) \cup \{A_{\mathrm{new}}\}, \\
  \Omega(K_{x+1}) &\subseteq
     \Omega(M_{x+1}) \cap
     \big\{ s \mid \Theta(K_{x+1})(s) \le 0 \big\}, \\
  M_{x+1}         &\supset M_x,
\end{align*}
все остальные компоненты \(K_{x+1}\) (потоки, потенциалы, циклы,
континуальность) задаются общей структурной машиной.

Операторы рождения \emph{минимальны} и \emph{необратимы}:
любой ненулевой вклад \(A_{\mathrm{new}}\) образует новый континуум
\(K_{x+1}\); возврат к прошлой размерности потребовал бы уничтожения
\(\Omega(K_{x+1})\), что интерпретируется как смерть, а не обратимое
упрощение.

\subsection{Оператор взаимодействия \texorpdfstring{$E_{\mathrm{int}}$}{E\_int}}

Когда взаимодействуют два континуума \(K_a\) и \(K_b\), их общая динамика
описывается оператором взаимодействия
\[
  E_{\mathrm{int}} : (K_a, K_b, M)
    \longrightarrow (K_a', K_b', M'),
\]
где \(M'\) — обновлённое пространство внедрения для объединённой системы.

На структурном уровне \(E_{\mathrm{int}}\) модифицирует потенциалы, потоки,
пороги, циклы и, возможно, оси каждого континуума. Различаются несколько
режимов:

\begin{itemize}
  \item \textbf{Конкуренция.} \
        Потоки одного континуума препятствуют поддержанию циклов в другом.
        Обычно \(k_a\) и \(k_b\) не могут оба бесконечно расти.
  \item \textbf{Паразитизм.} \
        Один континуум «похищает» поддерживающие потоки у другого, повышая
        собственную континуальность за счёт хозяина.
  \item \textbf{Симбиоз.} \
        Связанные потоки повышают континуальность обоих континуумов; могут
        возникать общие циклы, стабилизирующие пару.
  \item \textbf{Слияние.} \
        Оси и потенциалы объединяются в новый континуум
        \(K_{\mathrm{fusion}}\) с объединённым пространством состояний и
        порогами. Формально слияние соответствует рождению континуума
        большей размерности через составной оператор рождения.
\end{itemize}

Само взаимодействие может инициировать переходы размерности, когда смешанные
различия создают новые оси, недоступные континуумам по отдельности, и когда
напряжение превышает соответствующие пороги.

\subsection{Альгебра операторов и ограничения}

Хотя OC не постулирует полную алгебру операторов, для внутренней
согласованности требуются несколько структурных ограничений:

\paragraph{Совместимость с пространствами внедрения.}
Все операторы должны уважать включение
\(A(K(t)) \subseteq A(M(t))\) и не могут порождать эволюцию вдоль осей,
отсутствующих в \(M\). Операторы рождения определяются только тогда, когда
\(M\) уже содержит соответствующие оси.

\paragraph{Монотонность размерности.}
Структурный оператор \(S\) и операторы рождения \(\Psi_{x\to x+1}\) должны
удовлетворять
\[
  \dim A(t+dt) \ge \dim A(t),
\]
со строгим неравенством только на переходах размерности.

\paragraph{Динамика с учётом порогов.}
Операторы \(F, G, H, Q, R\) должны сохранять неравенственную структуру,
накладываемую порогами, за исключением разрешённых пересечений критических или
порогов смерти. Нарушения \(\Theta_{\mathrm{exist}}\) интерпретируются как
события смерти.

\paragraph{Континуальность как индикатор жизнеспособности.}
Оператор \(U\) связывает все остальные компоненты; в частности:
\begin{itemize}
  \item если \(k(t) = 0\) и \(\Omega = \emptyset\), континуум мёртв, и
        последующие применения \(E\) не дают эффекта на этом уровне;
  \item если \(k(t) > 0\), тогда существует по крайней мере один
        поддерживающий цикл, а потоки должны удовлетворять балансовым условиям,
        кодируемым в \(F\) и \(Q\).
\end{itemize}

\paragraph{Кросс-уровневая согласованность.}
Для переходов \(K_x \to K_{x+1}\) семейство операторов для более высокого
уровня должно сводиться к семейству более низкого, если новая ось фиксирована и
дополнительные степени свободы заморожены. Это обеспечивает восстановление
низших уровней как частных случаев высших, сохраняя вертикальную согласованность
иерархии.

\subsection{Итог}

Семейство операторов \((F, G, H, Q, R, S, U, \Psi, E_{\mathrm{int}})\)
представляет динамическую опору онтологии континуумов. Вместо единственного
высокоспециализированного уравнения движения OC использует модульный набор
операторов, действующих на потоки, потенциалы, пороги, циклы, границы, оси и
континуальность. Их совместное действие управляет рождением, жизнью,
взаимодействием и смертью континуумов на всех уровнях \(K_0\)–\(K_{12}\).

Эта глава завершает формальный каркас Core~1.3.3, делая явной структуру динамики,
которая была лишь имплицитно задана в ранних черновиках, и подготавливает
основу для полноценных количественных моделей в последующих расширениях.
